\section{Limitations}
\label{sec:limitations}

\paragraph{Statistical and schedule scope.}
Only three training seeds are available, with large variation for several
methods. We test one three-stage pacing schedule, without Reverse Task or
Reverse Expert. Moreover, staged eligibility changes cumulative sample
exposure as well as order, so these factors are not independently controlled.

\paragraph{Difficulty construction.}
Task and Expert components are manually chosen, and Expert duration, path
length, and cumulative action magnitude are redundant. Target radius enters
Task difficulty but is absent from the policy observation. Learner difficulty
depends on the MLP, its OOF training protocol, per-step averaging, trajectory
length, and phase composition; alternative choices could change the ranking.

\paragraph{Task and model coverage.}
Only pick-and-place and stacking are studied with small state-based MLPs.
Stacking has near-floor success and may not have converged to a useful regime.
It also uses relative features, precision-phase balancing, and stricter
persistent-contact success, preventing clean attribution of cross-task changes
to task identity.

\paragraph{Simulation and demonstration source.}
All demonstrations come from deterministic scripted experts in PyBullet. We
do not test human data, sensor or actuation noise, other policy families, real
robots, or sim-to-real transfer. The conclusions therefore concern this
controlled simulation study.
